\section{Declarations}\label{sec:declarations}
This section describes who is involved in the study, how the study handles data, ethics and research integrity, how AI tools were used in preparing the literature review and this proposal, and how the work was divided between the authors.

Two contacts at the University of Twente are involved in this study. Dr.ir.\ G.A.\ Folkertsma (Geert Folkertsma) is the topic contact for this research direction, and dr.ir.\ J.T.W.\ Berendsen (Jorien Berendsen) is the contact for the EE/ROB topics. The topic and methods are agreed with both contacts before the study starts, and the proposal is revised if the protocol is amended.

The study uses existing public datasets, namely SemanticKITTI and a synthetic dataset such as SynLiDAR, collects no new data and involves no human participants. The dataset licences, which are non-commercial for the KITTI family, are checked and respected, and no data is redistributed. Because no human participants or new data are involved, the study is not expected to require review by the faculty ethics committee, and this will be confirmed with Geert Folkertsma before the study starts.

To ensure research integrity and reproducibility, the seeds, configurations, class mapping and code are versioned in the shared repository. The threshold $\delta$ and the analysis plan are fixed before the degraded runs are evaluated, and the real evaluation data is not used for tuning. Results are reported for all degradation levels and diagnostics, including those that contradict the hypotheses.

Generative AI tools were used in preparing both documents. For the literature review, only OpenAI ChatGPT was used, for academic writing style, rephrasing and help with the search for literature. For this research proposal, the authors mainly used Anthropic's Claude to help with the writing, rephrasing and formatting of the text. After formulating the research gap, the authors also used Claude for an extensive check of whether this gap is indeed open. For both documents, the authors verified the correctness of the content and the references, and no AI tools were used for the research itself.

The authors contributed to this proposal as follows. Both authors defined the research direction in consultation with Geert Folkertsma. Frans Laporte derived the research gap from his part of the literature review on the evaluation of synthetic LiDAR data, wrote the proposal and revised it after the reviewer's feedback. Florent Kegler reviewed the revised proposal in detail, and his review led to changes in the hypotheses, the statistical analysis and the calibration of the baseline. He also set up the shared \LaTeX{} template and repository and corrected the references. Both authors reviewed the complete proposal and approve the final version. The division of work in the planned study is described in Section~\ref{sec:schedule}.